\documentclass[11pt,a4paper]{article}

\usepackage[utf8]{inputenc}
\usepackage[T1]{fontenc}
\usepackage{lmodern}
\usepackage{amsmath,amssymb,amsfonts}
\usepackage{booktabs}
\usepackage{hyperref}
\usepackage{microtype}
\usepackage{geometry}
\usepackage{cite}
\geometry{margin=2.5cm}

\title{\textbf{TaxonVision-MLOps: Architectural Concept Design \\ End-to-End Species Identification Pipeline}}
\author{\textbf{Benjamin Förster} \\ \texttt{mail@benjamin-foerster.de} \\ TaxonVision-MLOps Project}
\date{\today}

\begin{document}

\maketitle

\begin{abstract}
This document specifies the architectural concept design for \textbf{TaxonVision-MLOps}, an automated, production-grade Machine Learning Operations platform for fine-grained biological species identification from citizen-science imagery (e.g., iNaturalist, GBIF). Grounded in the principle that a simple, versioned, monitored, deployed, and automatically retrainable pipeline is superior to an isolated notebook model, this design addresses four graduated operational levels: closed-chain ingestion under strict Creative Commons licensing, Pareto-optimal backbone benchmarking and zero-allocation ONNX Runtime inference, closed-loop telemetry and taxonomic mutation reconciliation, and rigorous statistical safety guarantees via distribution-free Split Conformal Prediction ($1 - \alpha$ coverage) combined with active learning triage.
\end{abstract}

\section{Introduction and Problem Domain}
Ecological citizen science platforms such as iNaturalist record hundreds of millions of biodiversity observations~\cite{vanhorn2018inaturalist}. Developing an automated identification service for this domain presents fundamental challenges:
\begin{itemize}
    \item \textbf{Long-Tailed Taxonomic Distributions:} Taxa frequencies adhere to extreme power-law distributions where thousands of rare species possess few reference images.
    \item \textbf{Licensing and Provenance Governance:} Only observations distributed under open licenses (CC0, CC-BY, CC-BY-NC) may be ingested. Static domain assets lacking open licenses must be rejected, and attribution metadata must be preserved throughout the pipeline.
    \item \textbf{High-Stakes Decision Uncertainty:} Raw softmax probabilities are notoriously overconfident and uncalibrated. Mission-critical deployment demands mathematically guaranteed confidence sets with automatic referral to human experts for ambiguous observations.
\end{itemize}

\section{Graduated Architectural Levels}

\subsection{Level 1: Restricted Perimeter and Closed Chain}
The initial baseline establishes an automated, end-to-end chain over a representative subset of well-documented taxa:
\begin{itemize}
    \item \textbf{Data Ingestion:} Streaming and caching open-access observation metadata and images from the iNaturalist Open Data AWS S3 bucket.
    \item \textbf{Feature Extraction:} Fixed pre-trained feature extractors (e.g., lightweight convolutional or transformer architectures via \texttt{timm}, BioCLIP~\cite{stevens2024bioclip}) with a modular classification head.
    \item \textbf{Licensing Enforcement:} Automated cryptographic validation ensuring no non-open license images enter the dataset artifact.
    \item \textbf{Unified Web Service:} Minimal FastAPI serving endpoint exposing REST predictions alongside an HTMX operator dashboard.
\end{itemize}

\subsection{Level 2: Scalability, Cost, and Pareto Optimization}
Scaling to hundreds or thousands of taxa introduces cost and latency trade-offs:
\begin{itemize}
    \item \textbf{Pareto Backbone Benchmarking:} Systematic evaluation of candidate backbones (DINOv3, BioCLIP-2, DINOv2, MobileNetV4) balancing Top-1 accuracy against $p95$ latency and computational cost, trained with Class-Balanced Loss~\cite{cui2019classbalanced}.
    \item \textbf{Inference Optimization:} Exporting models to ONNX Runtime with dynamic INT8 and FP16 quantization, enforcing zero-allocation hot paths.
    \item \textbf{Large Dataset Versioning:} DVC-managed content-addressable storage maintaining reproducible dataset snapshots.
\end{itemize}

\subsection{Level 3: Feedback Loops and Production Observability}
Continuous operations require robust monitoring and label reconciliation:
\begin{itemize}
    \item \textbf{Closed Feedback Loop:} Automated ingestion of community-validated identifications over time to measure real-world performance degradation.
    \item \textbf{Taxonomic Mutation Handling:} Linnaean taxonomy revisions (splits, lumps, reclassifications) are tracked and mapped via DarwinCore taxon identifiers.
    \item \textbf{Telemetry:} Exporting Prometheus metrics for request latency, prediction entropy, human referral ratios, and feature embedding drift.
\end{itemize}

\subsection{Level 4: Mathematical Robustness, Cold Start, and Active Learning}
The safety layer protects downstream users and prioritizes human expert effort:
\begin{itemize}
    \item \textbf{Out-of-Distribution (OOD) Detection:} Energy-based scoring filters photographs devoid of organisms, corrupted images, or uncatalogued taxa.
    \item \textbf{Few-Shot \& Cold Start Protocol:} Embedding nearest-neighbor matching for rare taxa with limited photographic records.
    \item \textbf{Active Learning Triage:} Hybrid query strategies (BADGE and CoreSet)~\cite{ren2021deepactive} rank ambiguous samples to maximize information gain per human annotation.
    \item \textbf{Low-Latency Explanations:} Grad-CAM heatmap generation constrained to sub-25ms response budgets.
\end{itemize}

\section{Mathematical Invariants: Conformal Prediction}
Rather than relying on raw softmax scores $\hat{f}(X)$, we apply Split Conformal Prediction~\cite{angelopoulos2021gentle} to construct prediction sets $\hat{C}(X_{n+1}) \subseteq \mathcal{Y}$ satisfying:
\begin{equation}
    P(Y_{n+1} \in \hat{C}(X_{n+1})) \ge 1 - \alpha
\end{equation}
for a user-specified significance level $\alpha \in (0, 1)$.

Given $n$ holdout calibration samples $(X_i, Y_i)_{i=1}^n$, we define the non-conformity score:
\begin{equation}
    s_i = 1 - \hat{f}(X_i)_{Y_i}
\end{equation}
The critical calibration threshold $\hat{q}$ is computed as the $\lceil (n+1)(1-\alpha) \rceil / n$ empirical quantile of $\{s_1, \ldots, s_n\}$.
For a test image $X_{\text{test}}$, the output prediction set is:
\begin{equation}
    \hat{C}(X_{\text{test}}) = \left\{ y \in \mathcal{Y} : 1 - \hat{f}(X_{\text{test}})_y \le \hat{q} \right\}
\end{equation}
If $|\hat{C}(X_{\text{test}})| = 0$ (indicating an atypical or OOD observation) or $|\hat{C}(X_{\text{test}})| > k_{\text{max}}$ (indicating high taxonomic ambiguity), the observation is automatically referred to human triage.

\section{Conclusion}
The TaxonVision-MLOps concept design bridges cutting-edge representation learning and rigorous statistical guarantees within an automated, reproducible MLOps lifecycle.

\bibliographystyle{plain}
\bibliography{references}

\end{document}
